% authority: science_draft
\documentclass[11pt]{article}

\usepackage[margin=1in]{geometry}
\usepackage{amsmath,amssymb,amsthm}
\usepackage{booktabs}
\usepackage{enumitem}
\usepackage{longtable}

\newtheorem{definition}{Definition}
\newtheorem{theorem}{Theorem}
\newtheorem{proposition}{Proposition}
\newtheorem{remark}{Remark}
\newtheorem{nonconclusion}{Non-conclusion}

\newcommand{\Device}{\mathfrak D}
\newcommand{\Control}{\mathfrak C}
\newcommand{\Hist}{\mathcal H}
\newcommand{\Tokens}{\mathcal A}
\newcommand{\Response}{R}
\newcommand{\Difference}{\Delta}
\newcommand{\Malformed}{\bot_D}
\newcommand{\Undetermined}{\mathsf{undetermined}_D}
\newcommand{\Inconsistent}{\mathsf{inconsistent}_D}

\title{P7-T03 Source Operational-Device Suite Candidate v1}
\author{Candidate Constructor Packet RT-20260728-002}
\date{2026-07-28}

\begin{document}
\maketitle

\section{Control status and decisive result}

This fused task-local science draft constructs
\[
  \mathrm{SourceOperationalDeviceSuiteCandidate}_{v1}
\]
as a \emph{draft/control, proposal-only} finite source-protocol suite. It
defines explicit preparation, intervention, finite-history, readout-token,
success, and failure data for the guarded roles
\[
  \mathsf{Clock}_{\mathrm{src}},\quad
  \mathsf{Rod}_{\mathrm{src}},\quad
  \mathsf{Signal}_{\mathrm{src}},\quad
  \mathsf{Detector}_{\mathrm{src}},\quad
  \mathsf{FreeFall}_{\mathrm{src}}.
\]
It supplies two independent finite operational systems, exact rational
token-response maps, an intervention-difference map, token-partition and
source-presentation-naturality theorems, a signal--detector consistency
contract, and explicit positive, no-event, malformed, underdetermined, and
cross-device-inconsistent controls.

The role names are syntactic labels. A \(\mathsf{tick}\) token is not proper
time or duration; an endpoint or hop token is not physical length; an arrival
token is not a causal signal; a trigger token is not empirical detection; and
an unlabelled-intervention baseline is not gravitational free fall. History
weights are not physical probabilities, and the finite horizon is not
physical time.

Current ontology does not derive the protocol fields and does not adopt this
candidate. Adoption remains
\texttt{blocked\_adoption\_open\_continuation}. The result does not establish
effective geometry, physical clocks or rods, causality, detector semantics,
free-fall behavior, universal coupling, stress energy, a matter action,
Einstein equations, exact-GR recovery, benchmark promotion, proof authority,
publication, completed derivation, a global no-go theorem, or future
source-extension impossibility.
In particular, current ontology does not derive these protocols, and no
common physical propagation follows from their formal token relations.

\section{Fixed proposal-only source basis}

Hold the P7-T01
\(\mathrm{SourceMatterIncidenceCandidate}_{v1}\) fixed. Its admissible
configurations are \(\Psi=(m,z)\), and its typed transition records are
\[
  T=(\Psi^-,J,w,\Pi,\Psi^+)
\]
with source-only \(\Pi\) and the exact kinematic conditions
\[
  q(\Psi^+)-q(\Psi^-)=\partial_1J,\qquad
  z^+-z^-=\partial_2w.
\]
These conditions preserve the inherited sector coordinates
\((Q,[z])\) along an admissible transition, but they do not select an event,
provide an operational meaning, or establish a physical conservation law.

Also hold the P7-T02
\(\mathrm{SourceMatterFiniteTransitionKernelCandidate}_{v1}\) fixed. A
finite control is
\[
  \Control=(\Omega,E,s,t,K_{\Control}),
\]
where each finite enabled row \(E_x\) is nonempty and
\[
  \sum_{T\in E_x}K_{\Control}(T\mid x)=1.
\]
For a bookkeeping state
\(\mu\in\mathbb Q_{\geq0}^{\Omega}\), P7-T02 defines a formal finite update.
Its coefficients and composition counter carry no physical probability,
rate, duration, causal, or trajectory meaning. P7-T03 adds protocol data on
top of this fixed proposal-only object; it does not promote the kernel.

\section{Finite histories and protocol type}

\begin{definition}[Finite source history]
For \(n\in\mathbb N_0\), an admitted history is
\[
  h=(x_0,T_1,x_1,\ldots,T_n,x_n)\in\Hist_n(\Control)
\]
with \(s(T_j)=x_{j-1}\), \(t(T_j)=x_j\), and \(T_j\in E_{x_{j-1}}\).
Given initial bookkeeping state \(\mu\), define
\begin{equation}
  W_{\mu}(h)=
  \mu(x_0)\prod_{j=1}^{n}K_{\Control}(T_j\mid x_{j-1}).
  \label{eq:history-weight}
\end{equation}
This is a finite rational history weight. It is not a physical path
probability, amplitude, occurrence frequency, or empirical ensemble weight.
\end{definition}

\begin{definition}[Finite source operational protocol]
A protocol is the tuple
\[
  \Device=(\Control,P,I,\Tokens,r,S,F).
\]
Here:
\begin{enumerate}[leftmargin=*]
\item \(P\) is a finite nonempty set of preparation records
  \(p=(\Control_p,\mu_p,n_p)\);
\item \(I\) is a finite set of source intervention labels containing a
  baseline label \(i_0\); an intervention selects only declared source
  control and preparation data;
\item \(\Tokens\) is a finite formal token alphabet;
\item \(r\) is a total, single-valued, source-only map from every admitted
  finite history at the selected horizon to exactly one token;
\item \(S\subseteq\Tokens\) is a declared success-token subset; and
\item \(F\) classifies a request as valid success, valid no-event, valid
  protocol failure, malformed, underdetermined, or cross-device inconsistent.
\end{enumerate}
No member of the tuple may be defined by a target atlas, target metric,
proper time, physical distance, causal cone, empirical detector calibration,
desired benchmark answer, validator verdict, or generated derivative.
\end{definition}

\begin{definition}[Token response and intervention difference]
For a preparation \(p\), intervention \(i\), horizon \(n\), and token
\(a\in\Tokens\), define
\begin{equation}
  \Response_{\Device,p,i,n}(a)
  =
  \sum_{\substack{h\in\Hist_n(\Control_{p,i})\\r_{p,i,n}(h)=a}}
  W_{\mu_{p,i}}(h).
  \label{eq:response}
\end{equation}
Relative to the baseline \(i_0\), define
\begin{equation}
  \Difference_{\Device,p,i,n}(a)
  =
  \Response_{\Device,p,i,n}(a)
  -
  \Response_{\Device,p,i_0,n}(a).
  \label{eq:difference}
\end{equation}
Both maps are formal source bookkeeping records. They are not empirical
response functions.
\end{definition}

\section{Exact finite response theorems}

\begin{theorem}[Finite token-response closure and partition]
\label{thm:partition}
Let \(\Device\) be a valid protocol over a finite row-normalized control.
Then:
\begin{enumerate}[leftmargin=*]
\item every
  \(\Response_{\Device,p,i,n}(a)\in\mathbb Q_{\geq0}\);
\item
  \[
    \sum_{a\in\Tokens}\Response_{\Device,p,i,n}(a)
    =
    \sum_{x\in\Omega_{p,i}}\mu_{p,i}(x);
  \]
\item for a unit-bookkeeping-mass preparation, the token response partitions
  one; and
\item if baseline and intervention preparations have equal total
  bookkeeping mass, then
  \[
    \sum_{a\in\Tokens}\Difference_{\Device,p,i,n}(a)=0.
  \]
\end{enumerate}
\end{theorem}

\begin{proof}
Every history set is finite, and each factor in
Equation~\eqref{eq:history-weight} is a nonnegative rational, proving the
first claim. Totality and single-valuedness of \(r\) partition the history
set by tokens, so summing Equation~\eqref{eq:response} over \(a\) sums every
history weight exactly once. Repeated row normalization collapses the sum
over all length-\(n\) continuations to the initial bookkeeping mass. The unit
case follows. Subtracting the equal-mass baseline identity from the
intervention identity proves the zero-sum statement.
\end{proof}

\begin{nonconclusion}
The value one is a bookkeeping normalization. The theorem is not a
probability theorem, physical conservation law, completeness theorem for
observations, or empirical calibration statement.
\end{nonconclusion}

\begin{theorem}[Source-presentation naturality]
\label{thm:naturality}
Let \(g\) be a compatible P7-T01 presentation isomorphism transporting
P7-T02 controls row-bijectively. Suppose \(g\) also transports preparations,
interventions, and histories, and let
\(\alpha:\Tokens\to\Tokens'\) be a token bijection satisfying
\[
  r'(gh)=\alpha(r(h)).
\]
Then
\begin{equation}
  \Response_{g\Device,gp,gi,n}(\alpha(a))
  =
  \Response_{\Device,p,i,n}(a),
  \qquad
  \Difference_{g\Device,gp,gi,n}(\alpha(a))
  =
  \Difference_{\Device,p,i,n}(a).
  \label{eq:naturality}
\end{equation}
\end{theorem}

\begin{proof}
P7-T02 row-bijective transport preserves every kernel coefficient. Hence the
induced history bijection preserves the products in
Equation~\eqref{eq:history-weight}. The readout compatibility maps exactly
the histories in the \(a\)-fiber to those in the \(\alpha(a)\)-fiber.
Changing variables in the finite response sum proves the first equality;
subtracting transported baselines proves the second.
\end{proof}

This naturality is presentation-relative source bookkeeping. It is not
target covariance, diffeomorphism invariance, physical gauge symmetry, or
observational equivalence.

\section{Five guarded source protocol roles}

\subsection{\(\mathsf{Clock}_{\mathrm{src}}\)}

Take two source configurations \(c_0,c_1\) with rows
\[
 E_{c_0}=\{\mathbf1_{c_0},f_{01}\},\qquad
 E_{c_1}=\{\mathbf1_{c_1},f_{10}\},
\]
each coefficient being \(1/2\). Prepare
\(\mu_0=\delta_{c_0}\) at formal horizon two. Let the token be
\(\mathsf{tick}_1\) exactly for the transition word
\((f_{01},f_{10})\), and \(\mathsf{idle}_2\) otherwise. Then
\[
  R(\mathsf{tick}_1)=\frac14,\qquad
  R(\mathsf{idle}_2)=\frac34.
\]
Longer protocols may count declared completed cycle words and compare the
resulting integers. Neither the cycle count nor the fraction \(1/4\) is a
duration, proper time, clock rate, or frequency.

\subsection{\(\mathsf{Rod}_{\mathrm{src}}\)}

Take the ordered source-address chain
\[
  r_0\prec r_1\prec r_2
\]
with identity-plus-forward rows at \(r_0,r_1\) and an identity row at
\(r_2\). From \(\delta_{r_0}\) at horizon two, return
\(\mathsf{endpoint}_2\) exactly when the final address is \(r_2\).
The only successful word contains both forward hops, so
\[
  R(\mathsf{endpoint}_2)=\frac14.
\]
The declared order, endpoint index, and hop count are source-address data.
They are not length, distance, rigidity, a spatial coordinate, or a physical
rod standard.

\subsection{\(\mathsf{Signal}_{\mathrm{src}}\)}

Use a disjoint three-address source chain
\[
  s_0\longrightarrow s_1\longrightarrow s_2
\]
with the same identity-plus-forward row convention. The preparation carries
the token \(\mathsf{emitted}_0\). At horizon two, read
\(\mathsf{arrived}_2\) exactly at final address \(s_2\). Thus
\[
  R_{\mathrm{sig}}(\mathsf{arrived}_2)=\frac14,
\]
and two is the least formal horizon with nonzero arrival-token response.
This is not a light signal, causal arrival, propagation speed, cone, or
physical travel time.

\subsection{\(\mathsf{Detector}_{\mathrm{src}}\)}

Define source addresses
\[
  q_0=(s_0,d_{\mathrm{ready}}),\quad
  q_1=(s_1,d_{\mathrm{ready}}),\quad
  q_2=(s_2,d_{\mathrm{triggered}})
\]
with identity-plus-forward rows along \(q_0\to q_1\to q_2\). Return
\(\mathsf{triggered}_2\) exactly at \(q_2\). The declared history bijection
\[
  s_j\longleftrightarrow q_j
\]
and corresponding transition bijection imply
\[
  R_{\mathrm{sig}}(\mathsf{arrived}_2)
  =
  R_{\mathrm{det}}(\mathsf{triggered}_2)
  =
  \frac14.
\]
This is equality of two formal token fibers. It does not establish empirical
detection, measurement, observation, instrument response, or detector
universality.

\subsection{\(\mathsf{FreeFall}_{\mathrm{src}}\)}

Let the baseline label \(i_0\) mean only ``no declared source intervention.''
On configurations \(f_0,f_1\), take baseline row
\[
 E^{i_0}_{f_0}=\{\mathbf1_{f_0},m_{01}\}
\]
and an intervention \(i_{\mathrm{block}}\) with row
\[
 E^{i_{\mathrm{block}}}_{f_0}=\{\mathbf1_{f_0}\}.
\]
At horizon one from \(\delta_{f_0}\), return the final-address token. Then
\[
\begin{array}{c|cc}
 &\mathsf{address}_{f_0}&\mathsf{address}_{f_1}\\
\hline
i_0&1/2&1/2\\
i_{\mathrm{block}}&1&0\\
\Delta&1/2&-1/2
\end{array}
\]
and the difference sums to zero. ``No declared intervention'' is a syntactic
baseline, not force freedom. The contrast is not gravitational free fall,
geodesic motion, inertia, or an equivalence-principle test.

\section{Two independent systems and cross-device controls}

Declare the clock component \(S_{\mathrm{clock}}\) and signal component
\(S_{\mathrm{signal}}\) disjoint, and use their exact Cartesian product
control and product preparation. P7-T02 factorization gives
\[
  R_{\mathrm{joint}}
  (\mathsf{tick}_1,\mathsf{arrived}_2)
  =
  R_{\mathrm{clock}}(\mathsf{tick}_1)
  R_{\mathrm{signal}}(\mathsf{arrived}_2)
  =
  \frac1{16}.
\]
This verifies at least two independent formal operational systems. It is not
empirical statistical independence, locality, separability, or absence of a
physical force.

The suite also imposes three candidate-level consistency contracts:
\begin{enumerate}[leftmargin=*]
\item the signal arrival and detector trigger responses agree under their
  declared history bijection;
\item the disjoint clock--signal response factors only when the joint control
  and preparation are exactly Cartesian; and
\item rod endpoint and signal arrival addresses may be compared only when a
  declared source-address bijection transports their token labels.
\end{enumerate}
These checks compare formal source records. They do not derive common
physical propagation, metric universality, or universal coupling.

\section{Failure and residual branches}

\begin{longtable}{p{0.23\textwidth}p{0.69\textwidth}}
\toprule
Branch & Exact result\\
\midrule
\endhead
Valid positive &
All tuple fields are typed and at least one declared success token has
nonzero bookkeeping response.\\
Valid no-event &
All fields are typed, and an identity-only finite history returns a declared
no-event token. This is not a physical vacuum or stasis result.\\
Valid protocol failure &
All fields are typed, and the total readout returns a declared failure token.\\
Malformed &
\(\Malformed\) if a preparation or transition is untyped, the readout is
partial or multivalued, a token lies outside the alphabet, or any target
geometry or empirical calibration is used as a source premise.\\
Underdetermined &
\(\Undetermined\) when the source control is typed but preparation,
intervention, alphabet, readout, success, or failure data are absent. This is
not a no-go theorem.\\
Cross-device inconsistent &
\(\Inconsistent\) when a declared history bijection exists but paired token
fibers or response values disagree. This is a candidate-relative consistency
failure, not empirical falsification.\\
\bottomrule
\end{longtable}

For example, if the detector returns \(\mathsf{triggered}_2\) at either
\(q_1\) or \(q_2\), its response becomes \(3/4\) while the paired signal
arrival response stays \(1/4\). The declared contract therefore returns
\(\Inconsistent\). This local inconsistency does not prove that detector
semantics are impossible; a different source-only protocol remains open.

\section{Assumption and authority ledger}

The genuinely new proposal inputs are the finite preparations, intervention
labels, token alphabets, total readout maps, guarded role-name convention,
transport compatibility, and cross-device history bijections. From those
inputs and the fixed P7-T02 kernel, the candidate derives finite rational
responses, token partitions, zero-sum intervention differences,
source-presentation naturality, formal token-response equalities, and exact
Cartesian factorization.

No target metric, proper time, physical distance, causal cone, empirical
detector, force, coupling constant, dimensional scale, stress-energy object,
matter action, Einstein equation, benchmark behavior, validator status, or
generated derivative is an input.

The completion category for P7-T03 is
\texttt{constructed\_candidate}: the requested formal devices, at least two
independent systems, readout maps, interventions, success conditions, and
failure controls are explicit. This narrows only the proposal-level
operational-device sub-burden. It does not turn formal source tokens into
physical observables or adopt any primitive.

\section{Open burdens and next legal step}

Current ontology still does not derive the preparation, intervention, or
readout primitives. Physical interpretation and canonical adoption remain
open and protected. A later P7-T04 packet may ask whether multiple completed
formal source systems share a common source propagation structure, but only
after this packet is checkpointed. Even a common formal structure would not
by itself establish a lawful effective metric, universal coupling, physical
causality, or exact-GR recovery.

\end{document}
